\documentclass{article}

 
\usepackage[english]{babel}
 
\usepackage[letterpaper,top=2cm,bottom=2cm,left=3cm,right=3cm,marginparwidth=1.75cm]{geometry}
 
\usepackage{amsmath}
\usepackage{graphicx}
\usepackage[colorlinks=true, allcolors=black]{hyperref}

 
\usepackage{enumitem}
\setlist[enumerate,1]{label=\arabic*., itemsep=0.5em, parsep=0pt, topsep=0.5em}
\setlist[enumerate,2]{label=\alph*), itemsep=0.2em, parsep=0pt, topsep=0.2em, leftmargin=1.5em}
% --------------------------------------------

\title{Oncology Trial PI LLM Adoption Guide}
\author{Kevin Kawchak \\
\href{https://orcid.org/0009-0007-5457-8667}{0009-0007-5457-8667} \\
CEO ChemicalQDevice \\
\href{https://doi.org/10.5281/zenodo.20843290}{10.5281/zenodo.20843290}\\
\href{mailto:kevink@chemicalqdevice.com}{kevink@chemicalqdevice.com}}

\date{June 25, 2026} 


\begin{document}
\maketitle

\vspace{-0.75cm}
% Deadline: End of Day 6/24/26 (PDT)\\

\begin{abstract}
Repository based LLMs enable the creation of large scale documents more efficiently due to increased file generation capabilities, less errors, and a greater capacity to process an author's prompt containing specific instructions.\ This process allows for adequate incorporation of source information and iterative refinements of document drafts; while using a lower number of prompts. Repository level LLMs excel in understanding document context when a finite number of appropriately sized, high quality inputs are uploaded to the author's repository.\ Files generated in sections are advantageous for enabling direct access in future project builds.\ Single prompt multi-project outputs with a sub-prompt
topics schedule informs LLM generation and execution of sub-prompts possible for creating machine readable diagrams and three sets of self learning and iteratively refined clinical trial paper files.\ Both clinical trial site and sponsor simulations can also be conducted at scale to observe patient, robotic, and AI interactions.\ LLM proficiency is possible for at least glioblastoma, pancreatic ductal adenocarcinoma, and lung adenocarcinoma.\ This practical guide builds from the HHS June 22, 2026 initiative to restore American leadership in clinical trials. \\
\end{abstract}

\vspace{-0.5cm}

\tableofcontents

\vspace{0.1cm}

\section{Repository Setup}

\begin{enumerate}
\item Inputs:\ Maintain files accessed by the LLM in a single repository directory containing corresponding subdirectories to increase AI processing efficiency, document quality, and length.  
\begin{enumerate}
\item 6/23/26:\ Phase 2 Daraxonrasib, LLM Robotic PDAC Whipple, Opus 4.8, 64 pp. \cite{2026Ph2Darax}
\item 5/28/26: National Platform for Physical AI Oncology Trials, Opus 4.6, 186 pp. \cite{2026NatPlatf}
\end{enumerate}
\item Outputs:\ Generate files in separate sections when possible in order for text and code to serve as properly sized chunks for future downstream processing.  
\item Instruct LLM to create a markdown of your prompt in the repository at time of prompt submission.\ With prompts readily available on GitHub, future prompts can adapt from prior instructions to reduce future prompt writing.  \href{https://github.com/kevinkawchak/law-physical-ai-trials/blob/main/review/prompts/prompt-narrative.md}{[URL]}\cite{2026HR9510FL}
\item Pre-chunk large files into smaller sections, while instructing AI to generate a corresponding README of section contents and correlations between sections in order for the next main LLM processing step to execute more optimally. 
\begin{enumerate}
    \item NIH Phase 2, 3 Clinical Trials Protocol chunked into 10 sections, README. \href{https://grants.nih.gov/policy-and-compliance/policy-topics/clinical-trials/protocol-template}{[URL]}\href{https://github.com/kevinkawchak/physical-ai-oncology-trials/tree/main/trial-protocol/nih-protocol}{[URL]}\cite{2024nihprot, 2026OnPremis}
\end{enumerate}
\item Forked open-source repositories serve as living files, adaptable by Claude Code to create new projects at scale within hours; with code verifiable at any future date.  
\item LLMs can perform routine repository maintenance as prompts are submitted.
\begin{enumerate}
    \item “Make sure the repository is fully up to date with this work regarding badges, content, and context.” \href{https://github.com/kevinkawchak/Clinical-AI-Demos/blob/main/national-repositories/build-national.md}{[URL]}
    \item “Update other relevant documentation such as project structures, the main README diagrams, repository structure, etc. Update the CHANGELOG.md (v4.1.0)." \href{https://github.com/kevinkawchak/physical-ai-oncology-trials/blob/main/trial-phase-2/prompts/prompt-protocol.md}{[URL]}\cite{2026Ph2Darax}
    \item “There must be a single last v0.1.0 update that provides the main directory CHANGELOG.md, releases.md, repository updates, (as inspired by kevinkawchak/cancer-automated, but with 3x as many relevant main readme badges)." \href{https://github.com/kevinkawchak/law-physical-ai-trials/blob/main/review/prompts/prompt-narrative.md}{[URL]}\cite{2026HR9510FL}
\end{enumerate}
\item Human-in-the-loop monitorability is improved by specifying that the LLM continually push files to GitHub throughout the entire generation.  
\begin{enumerate}
    \item “Create an auto-commit / auto-PR process in real-time that allows for the user to monitor branch progress without any user intervention." \href{https://github.com/kevinkawchak/physical-ai-oncology-trials/blob/main/trial-protocol/prompts/prompt-protocol.md}{[URL]}\cite{2026OnPremis}
    \item “For each sub-prompt: 1 commit is required for each of the following (main.tex, .sty, .bib, and README); and 1 commit is required for each of the narrative's .tex sections (1 .tex file per section) that each correspond to main.tex." \href{https://github.com/kevinkawchak/physical-ai-oncology-trials/blob/main/trial-protocol/prompts/prompt-protocol.md}{[URL]}\cite{2026OnPremis}
\end{enumerate}
\item Reduce the size of README files to under 15K tokens to avoid LLM errors.
\end{enumerate}





\section{LLM Limitations}

\begin{enumerate}
\item Do not rely on Claude Code internet searches for large projects, especially those involving tabular and quantitative data. Specifying single URLs for Claude to access can be problematic. Instead, use a single GitHub repository with quality data of appropriate size for mission critical work. 
\item Claude Code does not always have access to compilers such as pdflatex for projects at scale. Instructing the LLM to make paper formatting improvements each iteration may be less effective, as the AI cannot view compilations.
\item Total input size should optimally be under 5MB, with the majority of the data being text, markdown, code, and machine readable images. 
\item LLM per file size limits exist outside of large 1M token context windows. 
\begin{enumerate}
\item READMEs and other documentation have increased to $\geq$25K tokens in recent releases.
\end{enumerate}
\item Be aware of providing LLM instructions that processes high dimensional data that would require approximations, and therefore affect deterministic outputs.
\item Server errors such as 404 and less interpretable errors are more common if the input is too large and/or the problem is too complex to solve. 
\begin{enumerate}
    \item Errors are the primary vehicle for the LLM to instruct to the author that output quality may be reduced. Additional prompts can often be submitted in this scenario to complete the generation, but is not advised, and may yield less than favorable results. 
\end{enumerate}
\item LLMs are trained more proficiently in some areas of expertise than others, such as new or de-emphasized fields \cite{HowLLMsTrain}. Double check the manufacturer's website to verify if specific benchmarks are available that focus on a specific area \cite{Opus48}.
\item More fundamental areas such as biology and chemistry have experienced access issues in recent months due to security issues \cite{SiliconCanals}.
\item Resulting context can be unexpected if the author is unsure of their dataset, which impacts the author's ability to prompt effectively.
\end{enumerate}

 

\section{Prompt Proficiency}

\begin{enumerate}
\item Sub-prompting is an effective method employed by first providing the LLM with a sub-prompt topics schedule \href{https://github.com/kevinkawchak/physical-ai-oncology-trials/blob/main/trial-protocol/prompts/prompt-protocol.md}{[URL]}. The LLM generates each respective sub-prompt \href{https://github.com/kevinkawchak/physical-ai-oncology-trials/tree/main/trial-protocol/sub-prompts}{[URL]}, and executes sub-prompts sequentially as project-scale outputs accumulate, all under a single master prompt. 
\begin{enumerate}
    \item A draft document, full document, final document, and embedded mermaid diagrams each in LaTeX were generated using a single prompt.\ The prompt contained LLM instructions to first create, and later execute sub-prompts. \cite{2026Ph2Darax, 2026OnPremis, 2026HR9510FL, 2026ClinTrst, 2026CongVote} 
\end{enumerate}
\item To improve observability, instruct the LLM to auto-commit to GitHub in real time so the author can review files as they are generated.
\item Meta-Prompting:\ Submit a prompt to the LLM with the dataset to receive a longer, more detailed, and more effective prompt. \href{https://github.com/kevinkawchak/Clinical-AI-Demos/blob/main/national-repositories/build-national.md}{[URL]}
\begin{enumerate}
    \item Meta-prompt LLM to return the paper draft template that includes code/text file repository locations to allow the next LLM processing step more effective.
    \item Utilize the generated prompt with the same dataset to process the files and directories.
    \item This method permits the LLM to focus on a single large processing task instead of having to find the correct files and directories in order to increase output length and quality.
\end{enumerate}
\item Visualization prompt additions permit human-in-the loop interaction during AI generations. 
\begin{enumerate}
    \item “Provide three ASCII diagrams, each with their own perspective, for each commit."
    % \item 
\end{enumerate}
\item AI should perform the heavy lifting; however workflows having many large tasks should be simplified and processed sequentially. 
\item Avoid excessive software engineering by relying on the LLM to perform the processing steps; and submit a follow-up processing report prompt after the primary generation is finished.
\begin{enumerate}
    \item “Provide a comprehensive set of processing metrics such as diff, number and type of agents used, number and type of subagents used, number and types of tools, skills, hooks, processing time, etc. Incorporate this information into several formatted text tables and ASCII diagrams. If there is a standard way of providing the information, provide this first - and then perform the rest of the tasks."
\end{enumerate}
\item Avoid GitHub tasks that are easily performed, such as transferring files to reduce token usage.
\item Some prompts have been further adapted for Claude Code, and are based on prior publications between 2024-2025. \cite{kawchak20}
\end{enumerate}





\section{Project Proficiency}

\begin{enumerate}
\item LLMs were shown to be proficient in clinical trial site and sponsor simulations at scale. \cite{Darax4Sim} 
\begin{enumerate}
    \item The reader has the option for the LLM to process all simulation data and return results; or generate agent files into the repository for use in subsequent steps. \cite{Darax4Sim} 
    \item LLMs have also been used to simulate future and expedited surgeries \cite{2026GlioAI60, 2026SixtySec}, a multi-robot, multi-site response team \cite{2026ThrHum24}, and for setting up a VVUQ oncology trial test suite. \cite{2026VVUQOnco} 
\end{enumerate}
\item Due to variable LLM pdf compilation ability, some author formatting is required to obtain publication quality. 
\begin{enumerate}
    \item Table column widths will likely require modification due to column character differences. 
    \item Consistent figure-caption vertical spacing usually requires individual edits per figure.
    \item Occasional paper text artifacts “=0mu plus 3mu” can be easily removed to fix table of contents running off the bottom of the page.
    \item Heading characters occasionally need to be reduced to avoid table of contents overlap.
    \item Additional clearpage commands should be utilized to advance a section to the next page. 
    \item Working in .md and .txt throughout projects reduces AI usage rates. If necessary, convert to .html, .xml, and other forms of markup in later steps. 
\end{enumerate}
\item First prompt generations typically have the largest outputs. Start new conversations, decrease the number of additional prompts, and reduce unnecessary content to improve context quality.
\begin{enumerate}
    \item BibTeX entries within .bib are typically correct, but long links can extend off of the right side of the pdf.
    \item BibTeX fixes are sometimes required if DOI and/or DOI URLs don’t print in pdf.  
    \item URLs should still be verified manually, although accuracy has increased in recent models.
    \item Single word lines commonly need to be addressed by adding or removing context.
    \item LLMs typically reward novel and well defined instruction based on repository text and code over optimizing similar prompt outputs using duplicate and triplicate prompt submissions.
    \item Spelling mistakes and sentence errors are uncommon with Claude Code, but context can vary based on author input data and prompt. \cite{kawchak20}
\end{enumerate}
\end{enumerate}

 

 
\section{Figure \& Table Tips}

\begin{enumerate}
\item Machine readable diagrams such as ASCII and Mermaid-type should be generated in favor of .png, .jpg, .bmp, especially in the beginning phases of projects. 
\begin{enumerate}
    \item “20+ Mermaid figures must be high quality, comprehensive, professional and professionally colored (20+ commits) (diagrams must be changed and improved throughout the draft, full, final process below)." \href{https://github.com/kevinkawchak/physical-ai-oncology-trials/blob/main/trial-protocol/prompts/prompt-protocol.md}{[URL]}\cite{2026OnPremis}
    \item Many diagrams can be observed throughout a generation.\ Three ASCII diagrams and four data files were generated on-demand every simulated hour for 24 hours.  \href{https://github.com/kevinkawchak/physical-ai-oncology-trials/tree/main/new-trial}{[URL]}  
    \item Mermaid and ASCII diagrams require some editing, especially if the LLM does not have access to a pdf compiler to learn from prior diagram iterations. 
\end{enumerate}
\item Multiple png chart/diagram generations can be performed using a meta-prompt with frontier LLMs when performed separately. 
\begin{enumerate}
    \item 10 or more high quality png images can be generated at a time by first instructing the LLM to generate an image generation prompt using a dataset; then using the generated prompt with the same dataset. \cite{2026MobPanCa}\href{https://github.com/kevinkawchak/cancer-automated/tree/main/papers/VVUQ-01}{[URL]}
    \item Only include placeholders in AI generated documents for raster images to be added in the final author edits; whereas machine readable diagrams can be used for each publication step. 
\end{enumerate}
\item \textbf{Tables:} Where data is clean and available, repository based LLMs now generate tables at scale. 
\begin{enumerate}
    \item Prior author tables can be used as templates for the LLM to learn from, making sure to either remove prior data or specify that prior tables are for formatting purposes only. 
    \item Specify that each table needs to be the width of the body text.
    \item Data conversions from .txt to .md to .tex complete seamlessly, as instructed to the LLM.
    \item Markup such as .html and .xml contain additional formatting which may affect scalability. 
    \item Longer tables by the LLM occasionally need to be split by the author into separate pages.
    \item Limitations stem from if the LLM currently has access to a pdf compiler, or not. \item Additional author instructions for the LLM to iterate on table column widths are limited to LLM compiler acccess. 
    \item Instruct the LLM to use \texttt{\textbackslash raggedright\textbackslash arraybackslash} to reduce inter-word spacing.
    \item “Make sure text doesn't run off the right side of the page anywhere." 
    \item Additional instructions are available in the Dec. 2025 LaTeX table prompting \cite{LLMTablePrompt}, Dec. 2025 Code scaling \cite{AIGenCode}, and Mar. 2026 Document scaling guides. \cite{PresentCode}
\end{enumerate}
\end{enumerate}
 

\clearpage 


\bibliographystyle{unsrt}
\bibliography{sample}

 
\section{Acknowledgments}
The author acknowledges ChatGPT 5.5 (Thinking Extended) and Gemini 3.1 Pro for manual editing assistance.\ Claude Code Opus and Fable models served as the main models specified in references. This v1.0 guide was assembled by Kevin Kawchak. 


\section{Ethical Disclosures}
This is an independent adoption guide. No real patients were enrolled, no protected
health information was used, and no IRB approval was sought or obtained in author references. This guide is not an active IND
or IDE, is not an approved protocol, and is not medical or regulatory advice; it is not endorsed by the FDA, NIH, HHS, an IRB, ICH, or any sponsor.


\section{Cite This Article}
Kawchak K. Oncology Trial PI LLM Adoption Guide. Zenodo.
2026; \href{https://doi.org/10.5281/zenodo.20843290}{10.5281/zenodo.20843290}.

\end{document}